\documentclass[conference]{assets/latex/IEEEtran}
\IEEEoverridecommandlockouts
% The preceding line is only needed to identify funding in the first footnote. If that is unneeded, please comment it out.
%Template version as of 6/27/2024

\usepackage{amsmath,amssymb,amsfonts}
\usepackage{algorithmic}
\usepackage{graphicx}
\usepackage{textcomp}
\usepackage{xcolor}
\usepackage[backend=biber,style=numeric]{biblatex}
\addbibresource{assets/latex/report.bib} 

 \graphicspath{{assets/latex/}}

\begin{document}

\title{CS6524 Computer Vision Lab. 2\\ Camera Calibration and Stereo Vision}

\author{
\IEEEauthorblockN{Andrew Jaffray}
\IEEEauthorblockA{23382163 \\ University of Limerick}
\and
\IEEEauthorblockN{Darragh Luby}
\IEEEauthorblockA{23359048 \\ University of Limerick}
\and
\IEEEauthorblockN{Patrick Walsh}
\IEEEauthorblockA{23372176 \\ University of Limerick}
}

\maketitle

\begin{abstract}
We calibrated our three phone cameras using a printed checkerboard, then we used each camera to estimate
depth by filming a video while moving sideways and treating two frames as a stereo pair. 
This calibration worked well, with errors under 0.25\,px for all three phones, and the depth maps show
near and far objects in the right places. The main problem is that the distances in centimetres
depend on a guess of how fast the camera was moving.
\end{abstract}

% Each section lives in its own file (owner noted at the top) to avoid merge conflicts.
% Andrew
% ============================ OWNER: Andrew ============================
\section{Introduction}
% TODO(Andrew): 3-4 sentences. Goal of the lab: calibrate three phone cameras, then estimate dense
% depth (cm) from a single moving camera by treating frame pairs as a stereo rig.

\section{Data Acquisition}
% TODO(Andrew): phones used (01 iPhone, 02 Android, 03 ?), portrait 720x1280 @ 30 fps,
% focus/exposure locked, board 9x6 squares (8x5 inner corners, s = 2.5 cm).
% Pre-processing identical for every video: ffmpeg, half resolution -> 360x640, H.264, no audio.

% Patrick
% ============================ OWNER: Patrick ===========================
\section{Camera Calibration}
% TODO(Patrick): method -- frame selection (board found with fast check, sharpest frame by Laplacian
% variance per time chunk, 20 frames), findChessboardCorners + cornerSubPix (11x11 window),
% object points on Z = 0 in cm, calibrateCamera, drop-worst-frame loop while error > 0.5 px.

\begin{table}[t]
\caption{Calibrated intrinsics and distortion (360$\times$640 px)}
\label{tab:calib}
\centering
\setlength{\tabcolsep}{3pt}
\begin{tabular}{lrrr}
\hline
 & \textbf{Cam 01} & \textbf{Cam 02} & \textbf{Cam 03} \\
\hline
$f_x$ & 453.74 & 454.71 & 537.70 \\
$f_y$ & 452.67 & 454.87 & 539.31 \\
$c_x$ & 188.75 & 183.05 & 183.58 \\
$c_y$ & 318.00 & 319.94 & 330.83 \\
skew & 0 & 0 & 0 \\
\hline
$k_1$ & $-0.123$ & $0.303$ & $0.045$ \\
$k_2$ & $6.365$ & $-10.158$ & $-1.576$ \\
$k_3$ & $-115.29$ & $95.24$ & $10.16$ \\
$p_1$ & $0.0007$ & $-0.0010$ & $0.0111$ \\
$p_2$ & $0.0049$ & $0.0020$ & $-0.0047$ \\
\hline
Reproj. err. (px) & 0.219 & 0.189 & 0.139 \\
\hline
\end{tabular}
\end{table}

% TODO(Patrick): discussion -- all errors well under the 0.5 px target, but k2/k3 are huge and of
% opposite sign between phones: the board never reached the image corners, so the high-order radial
% terms are unconstrained (overfit). Consequence: undistorting frames hurt pose estimation, so the
% depth stage uses K only (see Sec.~\ref{sec:depth}).

% Darragh
% ============================ OWNER: Darragh ===========================
\section{Inter-frame Pose and Depth Estimation}
\label{sec:depth}
% TODO(Darragh): method -- SIFT + ratio test, findEssentialMat (5-point RANSAC) instead of F -> E,
% recoverPose, pairs (t, t+10), reject pairs with |t_x| < 0.9 or rotation > 5 deg, stereoRectify +
% remap, StereoSGBM (64 disparities, block 5), Z = f B / d, B = assumed 15 cm/s x 10/30 s = 5 cm.
% Bonus: back-projected point cloud of a mid-video frame, dropping invalid / out-of-range / depth-edge pixels.

% TODO(Darragh): figure(s) -- one depth frame per camera + CloudCompare screenshot (put PNGs in assets/latex/).

\subsection{Discussion and Failure Analysis}
% TODO(Darragh): scale ambiguity (||t|| = 1, cm only via assumed speed, depth jumps between pairs),
% consecutive frames too close (tested gaps 5/10/15), F-route gave 180 deg flips, black left strip
% (numDisparities), glass/blank walls in 01-01, moving people in 01-02, floor-dominated 02 scenes.

% Andrew
% ============================ OWNER: Andrew ============================
\section{Improvements With More Time}
\textbf{Real scale.} Our depths depend on guessing that the camera moved at 15\,cm/s. If the
checkerboard were placed in each scene, \texttt{solvePnP} could give the camera's real position in
both frames, so we could measure how far it moved instead of guessing.

\textbf{Better calibration videos.} Our $k_2$ and $k_3$ values are far too large because the board
never reached the corners of the image. Recording again with the board moved into every corner, and
fixing $k_3=0$, should give distortion values we could actually use.

\textbf{Using more frames.} Each depth map comes from just two frames, so it flickers. Combining
several pairs would make it steadier and fill gaps, such as the plain wall in 03-01, where only 89
of 143 pairs were usable.



\printbibliography

\end{document}
